\chapter{Para explicarlo en el coloquio}
\label{cap:s-explicar}

\section{En un minuto}
\begin{cita}
Quintana Roo está lleno de un lado y vacío del otro. Los cinco lugares del sur tienen el 12.3\,\% de la gente del estado,
pero reciben el 1.4\,\% de los pasajeros de avión, y Tulum recibe 15.5 veces más visitantes que las pirámides del sur.
Nuestra campaña, ``El sur tiene espacio'', invita a la gente al sur sin llenarlo: con datos oficiales medimos dónde hay
espacio hoy, pronosticamos cuándo conviene ir y repartimos el presupuesto con un modelo de optimización que nunca anuncia
un lugar en temporada alta. Cada semana, una torre de control pausa los anuncios si hay mal clima o llega más gente de la
esperada. Con \$187,500 en 9 meses esperamos 957 visitantes más, a \$196 cada uno, y las reglas de cuidado no cuestan ni un
visitante.
\end{cita}

\section{En cinco minutos: la historia en orden}
\begin{enumerate}
\item \textbf{El problema} (capítulo 1): sargazo, Tulum en crisis, el norte lleno; la cifra 12.3\,\% contra 1.4\,\%.
\item \textbf{Los datos} (capítulos 3 y 4): 8.1 millones de registros oficiales, tres cajas, Spark; lo que no existe se
  dice.
\item \textbf{Dónde y cuándo} (capítulos 5 a 7): el índice de presión (los cinco lugares tranquilos), la forma del año
  (enero 61\,\% arriba en la Ruta), el pronóstico con su rango y las tormentas (agosto, 13.9\,\%).
\item \textbf{Cuánto y cómo se cuida} (capítulos 8 y 9): 957 visitantes, las reglas que cuestan cero, la Torre con 48
  pausas sin perder dinero.
\item \textbf{La campaña} (capítulo 10): dos viajeras, la marca, los anuncios con su dato y los 10 indicadores.
\end{enumerate}

\section{Si solo se recuerdan cinco cifras}
\begin{enumerate}
\item \textbf{12.3\,\% contra 1.4\,\%}: la gente del sur contra los pasajeros de avión que recibe.
\item \textbf{15.5 veces}: Tulum recibe 15.5 veces más visitantes que las pirámides del sur.
\item \textbf{957 visitantes por \$187,500}: lo que trae la campaña en 9 meses, a \$196 cada uno.
\item \textbf{0 visitantes}: lo que cuestan las reglas de cuidado ambiental y de equidad.
\item \textbf{250 extranjeros}: los que llegaron en avión a Chetumal en 2025, contra 9.4 millones a Cancún.
\end{enumerate}

\section{Preguntas probables y respuestas cortas}
\begin{teprofe}
\emph{``¿De dónde sacaste los datos?''} De 15 fuentes oficiales o abiertas (INEGI, SECTUR, el gobierno del estado, la NOAA,
el INAH por medio de SECTUR), descargadas con programas y guardadas sin modificar, cada archivo con su huella digital. Nada
de Google Maps ni TripAdvisor.
\end{teprofe}

\begin{teprofe}
\emph{``¿Por qué Spark?''} Porque la arquitectura tiene que seguir funcionando cuando los datos no quepan en una computadora:
el mismo código procesa 6 millones o 600 millones de registros. Se usó para limpiar los 6.1 millones de negocios del país y
para reproducir 239 semanas como si llegaran en vivo.
\end{teprofe}

\begin{teprofe}
\emph{``¿Por qué no estimaste la ocupación del sur?''} Porque nadie la midió desde 2025. Un modelo daría un número que parece
dato pero es inventado. Preferimos medir el sur con lo que sí se mide (tren, frontera, zonas arqueológicas) y decir ``sin
dato oficial''.
\end{teprofe}

\begin{teprofe}
\emph{``¿Cómo sabes que tu pronóstico funciona?''} Lo probamos ``viajando al pasado'': en cada mes pronosticamos los 12
siguientes solo con lo que se sabía entonces y lo comparamos con lo que pasó. El rango del 90\,\% se cumplió 91 de cada 100
veces en la Bahía, 94 en Belice y 80 en la Ruta, y decimos por qué la Ruta falla.
\end{teprofe}

\begin{teprofe}
\emph{``¿Qué es lo óptimo?''} Traer la mayor cantidad de visitantes con el presupuesto, sin romper cuatro reglas: no anunciar
en temporada alta, no pasar la capacidad probada, al menos 15\,\% por lugar y máximo 70\,\% por canal. Y lo comprobamos: las
reglas de cuidado cuestan cero visitantes.
\end{teprofe}

\begin{teprofe}
\emph{``¿Qué pasa si la campaña funciona demasiado bien?''} La regla anti-colapso está programada: la Torre pausa el anuncio
si llega más gente de la esperada, y el plan garantiza que ningún mes con anuncio pase del 40\,\% de la capacidad probada.
\end{teprofe}

\begin{teprofe}
\emph{``¿Por qué aparece Claude?''} Usamos un asistente de programación con IA como herramienta. Las reglas, las decisiones y
la revisión son nuestras y están documentadas con fecha en 25 notas de decisión.
\end{teprofe}

\section{Palabras que conviene saber explicar}
\begin{center}
\small
\begin{tabularx}{\linewidth}{>{\raggedright\arraybackslash}p{0.28\linewidth}Y}
\encabezado \h{Palabra} & \h{Qué significa, en una frase} \\
Bronce, plata y oro & Las tres cajas: lo crudo, lo limpio y lo listo para usar. \\
\filagris Hueco & Un dato que no existe; se declara y no se rellena. \\
Índice de presión & Un número de 0 a 1 que resume qué tan presionado está un lugar frente a todo el estado. \\
\filagris Forma del año & Cuántas veces un mes normal recibe cada mes (1.61 = 61\,\% más). \\
Capacidad probada & El mes más lleno que un lugar ha recibido en su historia. \\
\filagris Origen móvil & Probar un modelo ``viajando al pasado'', sin dejarlo ver el futuro. \\
Rango del 90\,\% & Entre qué valores caerá el dato real 9 de cada 10 veces. \\
\filagris Poisson & La regla para contar eventos raros que llegan sin aviso, como las tormentas. \\
Markov & Un modelo donde la siguiente semana depende solo de esta semana. \\
\filagris Monte Carlo & Imaginar miles de futuros posibles con azar para ver qué tan malo puede ser uno malo. \\
Simplex & El método que camina de esquina en esquina hasta encontrar el mejor reparto. \\
\filagris Precio sombra & Cuánto mejoraría el resultado si una regla se aflojara un poco. \\
Isolation Forest & Un método que detecta lo raro midiendo qué tan rápido se puede separar del resto. \\
\filagris Riesgo relativo & Cuántas veces más seguido pasa algo en un grupo que en general (ruido: 2.76). \\
Buyer persona & El retrato con datos de un cliente típico. \\
\filagris Spark Streaming & Procesar datos conforme llegan, en lotes pequeños y en orden. \\
\bottomrule
\end{tabularx}
\normalsize
\end{center}
